\documentclass[10pt,a4paper]{amsart}
\usepackage[margin=19mm]{geometry}
\usepackage{amsmath,amssymb,mathtools}
\usepackage[scr=rsfs,cal=euler]{mathalfa}
\usepackage{xfrac}
\usepackage[unicode,hidelinks,bookmarks=false]{hyperref}
\newcommand{\ie}{i.e.\ }
\newcommand{\cf}{cf.\ }
\newcommand{\resp}{respectively\ }
\newcommand{\Subsection}{\subsection}
\providecommand{\nobreakdash}{\nobreak\hbox{-}}
\setlength{\emergencystretch}{2em}
\multlinegap0pt
\usepackage{amssymb}
\usepackage{tikz}
\usepackage{algorithm}\usepackage{algpseudocode}
\newtheorem{proposition}{Proposition}
\title{\textbf{Frobenius–Galois expansions of substructural logics: Algebraization, Kalman equivalence and positive-cone semantics}}
\author{Juntao Wang$^{a, *}$, Jieqiong Shi$^{b}$, Mei Wang$^{c}$}
\date{Source: arXiv:2609.09529v1 (CC BY 4.0)}
\begin{document}
\maketitle

\noindent\emph{Source and licence.} \textbf{Frobenius–Galois expansions of substructural logics: Algebraization, Kalman equivalence and positive-cone semantics}. Version arXiv:2609.09529v1 (CC BY 4.0), distributed by arXiv under the Creative Commons Attribution 4.0 International licence, http://creativecommons.org/licenses/by/4.0/, as recorded in the arXiv licence metadata for this exact version and shown on the abstract page https://arxiv.org/abs/2609.09529v1. Full licence notice: \texttt{licenses/arxiv-2609.09529-CC-BY-4.0.txt}.\par\smallskip

\noindent\textit{Editorial scope.} Counted unit: the article's proposition prop:FGC-monotonicity (source0.tex lines 521--541). The article's complete original proof of it is delivered with it (source0.tex lines 543--588, one \texttt{proof} environment of 46 lines). No mathematics is written, repaired or re-derived: the statement and the proof are the article's own lines, in the article's own notation, and no retained line is substituted or re-flowed.  No further source-local context is needed: every cross-reference of every retained line resolves inside the two delivered spans.  Nothing was omitted from any retained span, so the package carries no omission record.  No retained line contains a citation, so the wrapper carries no bibliography and no bibliography entry.  No retained line contains a figure environment, an image-inclusion command or drawing code, so no image is delivered and the package declares no asset dependency.  Editorial formatting only, in addition: an AMS \texttt{amsart} preamble that restates the article's own declarations and macros used by the retained lines, namely the environments \texttt{proposition} on separate counters, together with the standard packages those lines need.  The retained environments print in this document as Proposition 1 in order of appearance, whereas the article prints the same items with the numbering of its own full document.  The complete native source of the article as distributed in the arXiv e-print for this exact version is bundled unchanged as \texttt{sources/209824/source0.tex}.
\begin{proposition}\label{prop:FGC-monotonicity}
Let $T$ be a theory of $\mathbf{FL}^{\mathrm{FGC}}_{ew}$ and
$\varphi,\psi\in\mathrm{Fm}_{\mathcal L_{\mathrm{FGC}}}$.

 (1)~If
 $T\vdash_{\mathbf{FL}^{\mathrm{FGC}}_{ew}}\varphi\Rightarrow\psi$, then
 $T\vdash_{\mathbf{FL}^{\mathrm{FGC}}_{ew}}
   \triangledown\varphi\Rightarrow\triangledown\psi
 ~\text{and}~
 T\vdash_{\mathbf{FL}^{\mathrm{FGC}}_{ew}}
   \blacktriangle\varphi\Rightarrow\blacktriangle\psi$,
 
 
 (2)~If
 $T\vdash_{\mathbf{FL}^{\mathrm{FGC}}_{ew}}
   \varphi\Leftrightarrow\psi$, then $T\vdash_{\mathbf{FL}^{\mathrm{FGC}}_{ew}}
   \triangledown\varphi\Leftrightarrow\triangledown\psi
 ~\text{and}~
 T\vdash_{\mathbf{FL}^{\mathrm{FGC}}_{ew}}
   \blacktriangle\varphi\Leftrightarrow\blacktriangle\psi$.
\end{proposition}

\begin{proof}
The reflexivity theorem
$\triangledown\varphi\Rightarrow\triangledown\varphi$, followed by
$(\mathrm{GC}_2)$ with $\varphi:=\triangledown\varphi$ and
$\psi:=\varphi$, gives
\[
 T\vdash_{\mathbf{FL}^{\mathrm{FGC}}_{ew}}
 \blacktriangle\triangledown\varphi\Rightarrow\varphi.
\]

If $T\vdash_{\mathbf{FL}^{\mathrm{FGC}}_{ew}}\varphi\Rightarrow\psi$, then by the transitivity of implication in $\mathbf{FL}_{ew}$,
\[
 T\vdash_{\mathbf{FL}^{\mathrm{FGC}}_{ew}}
 \blacktriangle\triangledown\varphi\Rightarrow\psi.
\]

Applying $(\mathrm{GC}_1)$ gives
\[
 T\vdash_{\mathbf{FL}^{\mathrm{FGC}}_{ew}}
 \triangledown\varphi\Rightarrow\triangledown\psi.
\]

For the second unary connective, start with the reflexivity theorem
$\blacktriangle\psi\Rightarrow\blacktriangle\psi$.  Applying $(\mathrm{GC}_1)$, with
$\varphi:=\psi$ and $\psi:=\blacktriangle\psi$, gives
\[
 T\vdash_{\mathbf{FL}^{\mathrm{FGC}}_{ew}}
 \psi\Rightarrow\triangledown\blacktriangle\psi.
\]

Combining this formula with $T\vdash_{\mathbf{FL}^{\mathrm{FGC}}_{ew}}\varphi\varphi\Rightarrow\psi$ yields
\[
 T\vdash_{\mathbf{FL}^{\mathrm{FGC}}_{ew}}
 \varphi\Rightarrow\triangledown\blacktriangle\psi.
\]

Applying $(\mathrm{GC}_2)$ now gives
\[
 T\vdash_{\mathbf{FL}^{\mathrm{FGC}}_{ew}}
 \blacktriangle\varphi\Rightarrow\blacktriangle\psi.
\]

(2) follows by applying (1) to the two
implications occurring in the definition of
$\varphi\Leftrightarrow\psi$.
\end{proof}

\end{document}
